\documentclass[11pt]{amsart}

% Encoding and fonts
\usepackage[T1]{fontenc}
\usepackage[utf8]{inputenc}
\usepackage{lmodern}

% Math packages
\usepackage{amsmath}
\usepackage{amssymb}
\usepackage{amsfonts}
\usepackage{amsthm}
\usepackage{mathtools}

% Layout and utilities
\usepackage{geometry}
\usepackage{enumitem}
\usepackage{microtype}
\geometry{margin=1in}

\numberwithin{equation}{section}
\setcounter{tocdepth}{3}
\setcounter{secnumdepth}{3}

% Theorem-like environments
\theoremstyle{plain}
\newtheorem{theorem}{Theorem}[section]
\newtheorem{proposition}[theorem]{Proposition}
\newtheorem{lemma}[theorem]{Lemma}
\newtheorem{corollary}[theorem]{Corollary}
\newtheorem{conjecture}[theorem]{Conjecture}

\theoremstyle{definition}
\newtheorem{definition}[theorem]{Definition}
\newtheorem{experiment}[theorem]{Experiment}
\newtheorem{heuristic}[theorem]{Heuristic}
\newtheorem{question}[theorem]{Question}

\theoremstyle{remark}
\newtheorem{remark}[theorem]{Remark}
\newtheorem{warning}[theorem]{Warning}

% Notation
\newcommand{\Mob}{\mu}
\newcommand{\Mert}{M}
\newcommand{\Qsf}{Q}
\newcommand{\Cat}{C}
\newcommand{\D}{\mathcal{D}}
\newcommand{\E}{\mathbb{E}}
\newcommand{\Pp}{\mathbb{P}}
\newcommand{\R}{\mathbb{R}}
\newcommand{\N}{\mathbb{N}}
\newcommand{\Fcat}{\mathcal{F}_{\mathrm{Cat}}}
\newcommand{\Icat}{\mathcal{I}_{\mathrm{Cat}}}
\newcommand{\eps}{\varepsilon}

% Hyperref must be last
\usepackage[hidelinks]{hyperref}

\title{Catalan Mertens}
\author{}
\date{}

\begin{document}

\maketitle
\tableofcontents
\bigskip

\begin{remark}
\label{rem:exploratory_nature}
This document is an exploratory mathematical notebook.  Its purpose is to isolate an exact Catalan structure inside the Mertens walk, to develop physical and probabilistic tools that act directly on that Catalan state space, and to identify a precise obstruction between ambient Catalan rarity and the deterministic arithmetic trajectory.  Rigorous statements are labelled as such; proposed mechanisms are labelled as heuristics, conjectures, or questions.
\end{remark}

\section{Introduction}
\label{sec:introduction}

Let
\begin{equation}
\Mert(x)=\sum_{n\le x}\Mob(n)
\label{eq:mertens_definition}
\end{equation}
be the Mertens function.  The classical square-root target is
\begin{equation}
\Mert(x)=O_{\eps}\!\left(x^{1/2+\eps}\right)
\qquad (\eps>0),
\label{eq:rh_mertens_target}
\end{equation}
which is equivalent to the Riemann hypothesis; see, for example, the standard analytic literature and the discussion in Bj\"orner's topological formulation \cite{Bjorner2011}.  The present notebook does not pursue an explicit-formula, zero-density, or Dirichlet-series proof strategy.  The target \eqref{eq:rh_mertens_target} is retained as the benchmark against which a different mechanism is measured.

The starting observation is elementary but exact.  Remove the zero values of the M\"obius function.  If
\begin{equation}
1=q_1<q_2<q_3<\cdots
\label{eq:squarefree_enumeration}
\end{equation}
are the squarefree positive integers and
\begin{equation}
X_m=\sum_{j\le m}\Mob(q_j),
\label{eq:compressed_walk}
\end{equation}
then every increment of $X_m$ is exactly $+1$ or $-1$, and
\begin{equation}
X_m=\Mert(q_m).
\label{eq:compressed_equals_mertens}
\end{equation}
Consequently, each interval between consecutive zeros of $X_m$ is a primitive lattice excursion.  If such an excursion has length $2n$, then, after fixing its sign, the number of possible excursion shapes is exactly the Catalan number
\begin{equation}
\Cat_{n-1}=\frac{1}{n}\binom{2n-2}{n-1}.
\label{eq:primitive_catalan_count_intro}
\end{equation}
Thus the real Mertens excursion is one state inside a finite Catalan state space.

The usefulness of this observation depends entirely on a second question.  Catalan combinatorics describes the ambient set of admissible excursion shapes, but the arithmetic trajectory is a deterministic selection from that set.  A rare Catalan shape remains possible unless the arithmetic selector itself satisfies an additional restriction.  The central problem of this notebook is therefore to identify a deterministic mechanism that controls the location of the Mertens trajectory inside Catalan state space.

The physical methods considered here are useful because they quantify geometry inside that state space.  Brownian excursion theory identifies the diffusive $\sqrt{L}$ scale; transfer matrices convert height restrictions into spectral gaps; path-integral arguments assign an action of order $H^2/L$ to an excursion of height $H$ and duration $L$; area-weighted Dyck paths provide genuine partition functions and free energies; Fredkin spin chains realize uniform superpositions of Dyck paths as quantum ground states; and TASEP provides a route from excursion statistics to non-equilibrium hydrodynamics and the Burgers equation \cite{Kaigh1976,DurrettIglehartMiller1977,KhorunzhiyMarckert2009,Kousha2012,OwczarekPrellberg2014,DerridaEnaudLebowitz2004,DellAnnaEtAl2016}.

The arithmetic side also has a nontrivial exact structure.  Bj\"orner associated to each cutoff $x$ a simplicial complex of squarefree integers and proved that the Mertens function is its reduced Euler characteristic \cite{Bjorner2011}.  In that representation, the compressed walk is the one-dimensional trace of a filtered, shifted simplicial complex.  This supplies a deterministic selector which is absent from a generic random walk.

The project therefore separates into three layers:
\begin{equation}
\text{Mertens filtration}
\longrightarrow
\text{Catalan excursion state}
\longrightarrow
\text{physical/entropic cost of its geometry}.
\label{eq:three_layer_program}
\end{equation}
The difficult arrow is the first-to-second selection law: why should the filtration select a non-exceptional region of the Catalan space?  The rest of the notebook develops several concrete formulations of this question.

\section{The selection wall}
\label{sec:selection_wall}

This section states the main obstruction at the outset.

\begin{definition}
\label{def:catalan_excursion_space}
Let $\D_n^{\mathrm{prim}}$ denote the set of positive primitive excursions of length $2n$: paths $D=(D_j)_{0\le j\le 2n}$ satisfying
\begin{equation}
D_0=D_{2n}=0,
\qquad
D_j>0 \quad (1\le j\le 2n-1),
\qquad
D_{j+1}-D_j\in\{-1,+1\}.
\label{eq:primitive_dyck_definition}
\end{equation}
\end{definition}

\begin{proposition}
\label{prop:primitive_catalan_cardinality}
For every $n\ge1$,
\begin{equation}
\left|\D_n^{\mathrm{prim}}\right|=\Cat_{n-1}.
\label{eq:primitive_catalan_cardinality}
\end{equation}
\end{proposition}

\begin{proof}
Delete the initial up-step and the final down-step and shift the remaining path down by one unit.  This is a bijection with ordinary Dyck paths of length $2n-2$, counted by $\Cat_{n-1}$.
\end{proof}

If a Mertens excursion has length $2n$, Proposition~\ref{prop:primitive_catalan_cardinality} places its shape in a set of size $\Cat_{n-1}$.  Under the uniform Catalan measure, high excursions occupy a small fraction of this set; the maximum of a uniform Dyck path is naturally of order $\sqrt n$, and exponential moments and moderate/large deviations are well studied \cite{KhorunzhiyMarckert2009,Kousha2012}.

\begin{warning}
\label{warn:rarity_not_selection}
Ambient rarity is not an arithmetic restriction.  A deterministic selector may choose a member of an exponentially small subset.  Therefore a statement of the form
\begin{equation}
\frac{\#\{D\in\D_n^{\mathrm{prim}}:\max D\ge H\}}
{\Cat_{n-1}}
\ll \exp\!\left(-cH^2/n\right)
\label{eq:ambient_rarity_only}
\end{equation}
becomes relevant to $\Mert$ only after one proves a property of the arithmetic selection rule.
\end{warning}

This warning is the main logical checkpoint of the project.  The physical theory can compute or estimate the cost of a region of Catalan space.  The arithmetical task is to show that the real trajectory cannot systematically occupy regions whose cost is too large.

\section{Exact Catalan decomposition of the compressed Mertens walk}
\label{sec:exact_catalan_decomposition}

\subsection{Compressed time}
\label{subsec:compressed_time}

Let
\begin{equation}
\Qsf(x)=\#\{n\le x:\Mob(n)\ne0\}
\label{eq:squarefree_count}
\end{equation}
be the squarefree counting function.  The standard asymptotic is
\begin{equation}
\Qsf(x)=\frac{6}{\pi^2}x+O(\sqrt{x}),
\label{eq:squarefree_density}
\end{equation}
so compressed time and ordinary integer time are linearly comparable.

\begin{proposition}
\label{prop:compression_identity}
For every $m\ge1$,
\begin{equation}
X_m=\Mert(q_m),
\label{eq:compression_identity}
\end{equation}
and for every real $x\ge1$,
\begin{equation}
\Mert(x)=X_{\Qsf(x)}.
\label{eq:decompression_identity}
\end{equation}
\end{proposition}

\begin{proof}
The terms omitted from \eqref{eq:compressed_walk} are exactly those for which $\Mob(n)=0$.
\end{proof}

\subsection{Return skeleton and factorization}
\label{subsec:return_skeleton}

Let
\begin{equation}
0=z_0<z_1<z_2<\cdots
\label{eq:return_times}
\end{equation}
be the successive return times of the compressed walk to zero, whenever they exist.  Put
\begin{equation}
L_r=z_r-z_{r-1}=2n_r.
\label{eq:excursion_lengths}
\end{equation}
Parity forces each complete excursion length to be even.

\begin{proposition}
\label{prop:catalan_product}
Fix a finite return skeleton $0=z_0<z_1<\cdots<z_R$ and fix the sign of each excursion.  The number of nearest-neighbour paths realizing exactly these complete excursions is
\begin{equation}
\mathcal{N}(z_1,\ldots,z_R)
=
\prod_{r=1}^{R}\Cat_{n_r-1}.
\label{eq:catalan_product}
\end{equation}
\end{proposition}

\begin{proof}
The excursion interiors are independent combinatorial choices once their durations and signs are specified.  Proposition~\ref{prop:primitive_catalan_cardinality} applies to each factor.
\end{proof}

Equation \eqref{eq:catalan_product} is an exact entropy associated with the zero-return skeleton:
\begin{equation}
S_{\mathrm{Cat}}(z_R)
=
\log \mathcal{N}(z_1,\ldots,z_R)
=
\sum_{r=1}^{R}\log \Cat_{n_r-1}.
\label{eq:catalan_skeleton_entropy}
\end{equation}
Using Stirling's formula,
\begin{equation}
\log \Cat_n
=
n\log 4-\frac{3}{2}\log n+O(1).
\label{eq:catalan_entropy_asymptotic}
\end{equation}
Thus each completed excursion carries an entropy deficit of order $\frac32\log n$ relative to unrestricted $\pm1$ words of comparable length.

\section{Open Catalan sectors: meanders and pointwise control}
\label{sec:meanders}

Complete excursions alone do not control a point lying inside a very long unfinished excursion.  The natural pointwise object is therefore a Catalan prefix, or Bernoulli meander.  This is also the correct continuum bridge to Brownian meander \cite{DurrettIglehartMiller1977,CsakiMohanty1981}.

\begin{definition}
\label{def:positive_meander}
For $t\ge1$ and $h>0$ with $t\equiv h\pmod 2$, let $\mathcal{M}_{t,h}$ be the set of paths $S_0=0$, $S_t=h$, with increments $\pm1$ and
\begin{equation}
S_j>0\qquad (1\le j\le t).
\label{eq:strict_meander_condition}
\end{equation}
Let $\mathcal{M}_t=\bigcup_h\mathcal{M}_{t,h}$.
\end{definition}

\begin{proposition}
\label{prop:ballot_meander_count}
The ballot theorem gives
\begin{equation}
\left|\mathcal{M}_{t,h}\right|
=
\frac{h}{t}
\binom{t}{(t+h)/2},
\label{eq:ballot_endpoint_count}
\end{equation}
and
\begin{equation}
\left|\mathcal{M}_t\right|
=
\binom{t-1}{\lfloor (t-1)/2\rfloor}.
\label{eq:total_meander_count}
\end{equation}
\end{proposition}

\begin{remark}
\label{rem:meander_as_open_catalan}
The meander is the open-boundary sector of the same Catalan geometry.  Closing a positive meander by a first return produces a primitive Dyck excursion, while cutting a primitive Dyck excursion before its terminal return produces a positive meander.  For pointwise estimates on $\Mert$, the meander formulation avoids any dependence on a future return time.
\end{remark}

For a compressed time $m$ let
\begin{equation}
a(m)=\max\{z_r:z_r\le m\},
\qquad
t(m)=m-a(m),
\qquad
h(m)=|X_m|.
\label{eq:last_zero_coordinates}
\end{equation}
Whenever $h(m)>0$, the segment from $a(m)$ to $m$ is an element of $\mathcal{M}_{t(m),h(m)}$ after reflection if necessary.

\begin{proposition}
\label{prop:meander_gaussian_tail}
There exist absolute constants $c,C>0$ such that, under the uniform measure on $\mathcal{M}_t$,
\begin{equation}
\Pp_t\{S_t\ge h\}
\le
C\left(1+\frac{h}{\sqrt t}\right)
\exp\!\left(-c\frac{h^2}{t}\right)
\label{eq:meander_endpoint_tail}
\end{equation}
for admissible $t,h$.
\end{proposition}

\begin{proof}
Equation \eqref{eq:ballot_endpoint_count} reduces the endpoint law to binomial coefficients.  Standard Chernoff bounds for the symmetric binomial distribution give Gaussian decay in $h^2/t$, while \eqref{eq:total_meander_count} is of order $2^t/\sqrt t$.  Summing the endpoint masses from $h$ upward gives \eqref{eq:meander_endpoint_tail} after adjustment of absolute constants.
\end{proof}

This proposition supplies a rigorous ambient energy scale at the current time, not only at completed returns.

\section{Catalan free energy and a sufficient selection principle}
\label{sec:catalan_free_energy}

\begin{definition}
\label{def:meander_selection_energy}
For the actual compressed Mertens path at a time $m$ with $h(m)>0$, define its open Catalan selection energy by
\begin{equation}
\Icat(m)
=
-\log
\Pp_{t(m)}\{S_{t(m)}\ge h(m)\},
\label{eq:selection_energy_definition}
\end{equation}
where $\Pp_t$ is the uniform measure on $\mathcal{M}_t$.
\end{definition}

This is a deterministic number attached to the actual path: probability is used only to measure the volume of the ambient Catalan/meander sector containing an endpoint at least as extreme as the observed one.

\begin{proposition}
\label{prop:selection_energy_lower_bound}
There are constants $c_1,c_2,c_3>0$ such that
\begin{equation}
\Icat(m)
\ge
c_1\frac{h(m)^2}{t(m)}
-c_2\log\!\left(2+\frac{h(m)}{\sqrt{t(m)}}\right)
-c_3.
\label{eq:selection_energy_lower_bound}
\end{equation}
\end{proposition}

\begin{proof}
Take minus logarithms in Proposition~\ref{prop:meander_gaussian_tail}.
\end{proof}

\begin{theorem}
\label{thm:selection_energy_implies_rh_scale}
Assume that for some constant $A>0$ and all sufficiently large compressed times $m$,
\begin{equation}
\Icat(m)\le A\log m.
\label{eq:selection_energy_hypothesis}
\end{equation}
Then
\begin{equation}
|X_m|=O\!\left(\sqrt{m\log m}\right),
\label{eq:compressed_sqrtlog_bound}
\end{equation}
and consequently
\begin{equation}
\Mert(x)=O\!\left(\sqrt{x\log x}\right).
\label{eq:mertens_sqrtlog_bound}
\end{equation}
In particular, \eqref{eq:rh_mertens_target} follows.
\end{theorem}

\begin{proof}
By Proposition~\ref{prop:selection_energy_lower_bound} and \eqref{eq:selection_energy_hypothesis}, the quantity $h(m)^2/t(m)$ is $O(\log m)$, up to a lower-order logarithmic correction.  Since $t(m)\le m$, this gives \eqref{eq:compressed_sqrtlog_bound}.  Equation \eqref{eq:decompression_identity} and the linear comparability \eqref{eq:squarefree_density} then give \eqref{eq:mertens_sqrtlog_bound}.
\end{proof}

\begin{warning}
\label{warn:selection_energy_is_the_wall}
The hypothesis \eqref{eq:selection_energy_hypothesis} is precisely the unresolved selector statement.  Theorem~\ref{thm:selection_energy_implies_rh_scale} is useful because it isolates the target in a finite combinatorial language: the arithmetic path must be prevented from entering meander sectors whose relative Catalan volume is smaller than a power of the current scale.
\end{warning}

\section{Brownian excursion and meander as continuum Catalan mechanics}
\label{sec:brownian_mechanics}

Kaigh proved an invariance principle for a recurrent random walk conditioned on a late first return to zero: after diffusive rescaling the path converges to Brownian excursion \cite{Kaigh1976}.  Durrett, Iglehart, and Miller established weak convergence results for Brownian meander and excursion in related conditional settings \cite{DurrettIglehartMiller1977}.  These results provide a continuum mechanics for the Catalan state space rather than an independent model imposed on the arithmetic data.

For an excursion of discrete duration $L$, the diffusive scaling is
\begin{equation}
y_L(s)=L^{-1/2}D_{\lfloor Ls\rfloor},
\qquad 0\le s\le1.
\label{eq:diffusive_excursion_scaling}
\end{equation}
The limiting excursion has order-one height, so the natural discrete height scale is
\begin{equation}
H\asymp \sqrt L.
\label{eq:brownian_height_scale}
\end{equation}
The exponential moments of the normalized Dyck maximum are uniformly controlled and converge to the corresponding Brownian-excursion quantity \cite{KhorunzhiyMarckert2009}.  Kousha further developed large and moderate deviation principles for the Dyck maximum \cite{Kousha2012}.

\subsection{A variational action}
\label{subsec:variational_action}

\begin{heuristic}
\label{heur:brownian_action}
This is a physical large-deviation heuristic.  For a continuum path $y:[0,L]\to\R_+$ with $y(0)=y(L)=0$, assign the Euclidean action
\begin{equation}
\mathcal{A}[y]
=
\frac{\kappa}{2}
\int_0^L \dot y(t)^2\,dt.
\label{eq:euclidean_action}
\end{equation}
Among paths that reach height $H$, the piecewise-linear triangular profile minimizes \eqref{eq:euclidean_action}.  Its action is
\begin{equation}
\mathcal{A}_{\min}(H,L)
=
2\kappa\frac{H^2}{L}.
\label{eq:triangular_action}
\end{equation}
Thus the physical rate variable is $H^2/L$, matching the Gaussian scale already present in the exact ballot estimate \eqref{eq:meander_endpoint_tail}.
\end{heuristic}

The importance of \eqref{eq:triangular_action} is conceptual: the physical energy scale is not being guessed from Mertens data.  It is the continuum expression of a rate already visible in Catalan/ballot combinatorics.

\section{Transfer matrices and the quantum Catalan picture}
\label{sec:transfer_quantum}

\subsection{Height confinement as a spectral problem}
\label{subsec:height_spectral}

Let $A_h$ be the adjacency matrix of the path graph on heights $0,1,\ldots,h$.  The number of Dyck paths of length $2n$ confined to this height interval is
\begin{equation}
Z_n(h)=(A_h^{2n})_{00}.
\label{eq:confined_partition_function}
\end{equation}
The eigenvalues are
\begin{equation}
\lambda_k(h)
=
2\cos\!\left(\frac{k\pi}{h+2}\right),
\qquad
1\le k\le h+1.
\label{eq:path_graph_eigenvalues}
\end{equation}
Consequently,
\begin{equation}
Z_n(h)
=
\frac{2}{h+2}
\sum_{k=1}^{h+1}
\sin^2\!\left(\frac{k\pi}{h+2}\right)
\lambda_k(h)^{2n}.
\label{eq:confined_spectral_sum}
\end{equation}
This spectral representation is standard in the analysis of Dyck-path maxima and underlies related deviation results \cite{Kousha2012}.

For large $h$,
\begin{equation}
\lambda_1(h)
=
2-\frac{\pi^2}{(h+2)^2}+O(h^{-4}),
\label{eq:ground_eigenvalue_expansion}
\end{equation}
so confinement below an atypically low ceiling has a free-energy scale of order
\begin{equation}
\Delta F_{\mathrm{conf}}
\asymp
\frac{n}{h^2}.
\label{eq:confinement_free_energy}
\end{equation}
This is complementary to the high-excursion action $H^2/n$ in \eqref{eq:triangular_action}.  The two scales are the spectral and variational faces of the same diffusive geometry.

\subsection{Fredkin spin chains}
\label{subsec:fredkin}

The Fredkin spin chain provides a genuine quantum realization of the Catalan ensemble.  In the spin-$1/2$ model, the ground state can be written as an equal-weight superposition of Dyck paths; related models and their correlation structure are developed in \cite{DellAnnaEtAl2016,SalbergerKorepin2017}.  Abstractly, if $\D_n$ is the Dyck basis,
\begin{equation}
|\Psi_n\rangle
=
\frac{1}{\sqrt{\Cat_n}}
\sum_{D\in\D_n}|D\rangle.
\label{eq:fredkin_catalan_state}
\end{equation}
For any subset $\mathcal{A}\subseteq\D_n$ with projector $\Pi_{\mathcal A}$,
\begin{equation}
\langle\Psi_n|\Pi_{\mathcal A}|\Psi_n\rangle
=
\frac{|\mathcal A|}{\Cat_n}.
\label{eq:sector_weight_quantum}
\end{equation}
Thus a Catalan free energy
\begin{equation}
\Fcat(\mathcal A)
=
-\log\frac{|\mathcal A|}{\Cat_n}
\label{eq:catalan_free_energy_sector}
\end{equation}
has a literal interpretation as the negative logarithm of a quantum-sector weight in an equal-amplitude Catalan state.

\begin{question}
\label{question:arithmetic_hamiltonian}
Can the arithmetic selector be represented by a local or quasi-local deformation of a Catalan transfer matrix or Fredkin-type Hamiltonian whose ground-state or low-energy sector contains the Mertens excursion?  A positive answer would turn the selector problem into a spectral-energy estimate rather than a statement of random sampling.
\end{question}

\section{Statistical mechanics of Dyck paths: pressure, area, and $q$-Catalan weights}
\label{sec:dyck_statistical_mechanics}

Dyck paths are established physical configuration spaces for directed polymers and vesicles.  Owczarek and Prellberg study area-weighted Dyck paths as a two-dimensional vesicle model and define pressure through finite-size free-energy differences \cite{OwczarekPrellberg2014}.  Let $A(D)$ be the area under a Dyck path and define
\begin{equation}
Z_n(q)
=
\sum_{D\in\D_n} q^{A(D)}.
\label{eq:q_catalan_partition}
\end{equation}
At $q=1$,
\begin{equation}
Z_n(1)=\Cat_n,
\label{eq:q_catalan_at_one}
\end{equation}
while $q\ne1$ biases the ensemble by area.  The asymptotics of area-weighted Dyck paths exhibit Airy-type scaling and have been analyzed in detail \cite{HaugPrellberg2015}.

For a Mertens excursion $D_r$, define its absolute area
\begin{equation}
A_r
=
\sum_{j=z_{r-1}}^{z_r}|X_j|.
\label{eq:mertens_excursion_area}
\end{equation}
The pair $(H_r,A_r)$ separates a sharp spike from a sustained displacement.  A physical route based on $q$-Catalan weights would seek an arithmetic inequality on the partition weight of the sector containing $(H_r,A_r)$.

\begin{heuristic}
\label{heur:entropic_pressure}
In a Catalan polymer interpretation, excluding a region of path space produces an entropic pressure equal to a free-energy difference.  An arithmetic excursion that remains at large height for a long interval can therefore be assigned a pressure or free-energy cost determined entirely by the number of Catalan paths compatible with that geometric constraint.  The open problem is to connect this cost to the deterministic arithmetic selector.
\end{heuristic}

\section{Non-equilibrium transport and the Burgers route}
\label{sec:tasep_burgers}

The asymmetric exclusion process provides a distinct physical use of Brownian excursions.  Derrida, Enaud, and Lebowitz showed that in the maximal-current phase of open TASEP, the non-Gaussian part of density fluctuations can be described in terms of Brownian excursion statistics \cite{DerridaEnaudLebowitz2004}.  At hydrodynamic scale, TASEP is governed by a conservation law of Burgers type,
\begin{equation}
\partial_t\rho
+
\partial_x\bigl(\rho(1-\rho)\bigr)
=0.
\label{eq:tasep_burgers}
\end{equation}
This creates a rigorous physical chain linking constrained lattice paths, excursion fluctuations, current transport, and non-linear hydrodynamics.

The present project does not identify the Mertens walk with TASEP.  The possible use is structural: construct a dynamics on Catalan states whose conserved current or entropy production encodes the arithmetic selector.  If the selector can be expressed as a constrained transport process, tools from hydrodynamic limits, entropy solutions, shocks, and large-deviation functionals become available.

\begin{question}
\label{question:selector_as_current}
Is there a local state variable on the filtered squarefree complex for which the parity imbalance $X_m$ is a conserved or approximately conserved current, and for which the admissible Catalan updates induce a non-equilibrium transport law?
\end{question}

A positive answer could replace a direct estimate on $\Mob(n)$ by an entropy-production inequality on the Catalan dynamics.  This is speculative, but the TASEP/Brownian-excursion/Burgers chain shows that such a passage from discrete constrained paths to hydrodynamic free energies is a real mechanism in statistical physics, not merely a metaphor.

\section{The deterministic arithmetic selector}
\label{sec:arithmetic_selector}

\subsection{Björner's filtered complex}
\label{subsec:bjorner_complex}

For each squarefree integer $k$, let $P(k)$ be the set of its prime divisors and define
\begin{equation}
\Delta_x
=
\{P(k): k\le x,\ k\text{ squarefree}\}.
\label{eq:bjorner_complex}
\end{equation}
This is an abstract simplicial complex because divisors of a squarefree integer correspond to subsets of its prime set.  Bj\"orner observed and exploited the exact identity
\begin{equation}
\Mert(x)
=
-\widetilde\chi(\Delta_x),
\label{eq:mertens_euler_characteristic}
\end{equation}
where $\widetilde\chi$ is the reduced Euler characteristic \cite{Bjorner2011}.  He further showed that these complexes are shifted and studied their Betti numbers, wedge-of-spheres homotopy type, and combinatorial shadow inequalities.  Bj\"orner explicitly noted that these results did not by themselves yield new control of the growth of $\Mert(x)$ \cite{Bjorner2011}.

The filtration
\begin{equation}
\Delta_1\subseteq\Delta_2\subseteq\Delta_3\subseteq\cdots
\label{eq:simplicial_filtration}
\end{equation}
adds a simplex precisely when the cutoff crosses a squarefree integer.  The change in Euler characteristic is therefore the sign of the corresponding M\"obius value.  The compressed Mertens walk is exactly the Euler-characteristic trace of this filtration.

\subsection{Subset-sum energy representation}
\label{subsec:subset_sum_energy}

A squarefree integer corresponds to a finite prime subset $S$, and its logarithm is additive:
\begin{equation}
E(S)
=
\log\prod_{p\in S}p
=
\sum_{p\in S}\log p.
\label{eq:subset_sum_energy}
\end{equation}
The sign is the parity
\begin{equation}
\Mob\!\left(\prod_{p\in S}p\right)
=
(-1)^{|S|}.
\label{eq:subset_parity_sign}
\end{equation}
Thus the selector can be stated exactly as follows: order all finite prime subsets by the additive energy \eqref{eq:subset_sum_energy}, then record the parity \eqref{eq:subset_parity_sign}.  Mertens is the cumulative parity imbalance.

For finitely many primes $p_1,\ldots,p_d$, let
\begin{equation}
w=(w_1,\ldots,w_d),
\qquad
w_j=\log p_j,
\label{eq:prime_weight_vector}
\end{equation}
and
\begin{equation}
E_w(S)=\sum_{j\in S}w_j.
\label{eq:weighted_subset_energy}
\end{equation}
The relative order of two subsets changes on a hyperplane
\begin{equation}
H_{S,T}
=
\left\{w:\langle \mathbf{1}_S-\mathbf{1}_T,w\rangle=0\right\}.
\label{eq:subset_order_hyperplane}
\end{equation}
Unique factorization implies that the prime-log vector does not lie on any hyperplane $H_{S,T}$ with $S\ne T$.

\begin{proposition}
\label{prop:subset_gap}
If $a$ and $b$ are distinct positive integers with $a,b\le X$, then
\begin{equation}
|\log a-\log b|
\ge
\log\!\left(1+\frac{1}{X}\right).
\label{eq:log_integer_gap}
\end{equation}
Consequently, the ordering of any fixed finite family of squarefree products $\le X$ is separated from its subset-sum collision hyperplanes by a positive deterministic margin.
\end{proposition}

\begin{proof}
Assume $a>b$.  Then $a\ge b+1$ and $b\le X$, so
\begin{equation}
\log(a/b)
\ge
\log(1+1/b)
\ge
\log(1+1/X).
\label{eq:log_gap_proof}
\end{equation}
\end{proof}

This hyperplane-arrangement formulation gives a concrete geometric meaning to ``selection'': a finite parity word, and hence a finite Catalan excursion, is stable on a chamber of weight space.  The arithmetic vector $(\log p_j)$ lies in one such chamber.

\begin{definition}
\label{def:selector_chamber}
For a finite family $\mathcal S$ of prime subsets, define the selector chamber of the prime-log vector to be the connected component containing $w=(\log p_j)$ of
\begin{equation}
\R_{>0}^{d}
\setminus
\bigcup_{S\ne T\in\mathcal S}H_{S,T}.
\label{eq:selector_chamber_definition}
\end{equation}
\end{definition}

The chamber solid angle or normalized volume is a possible deterministic analogue of statistical weight.  A useful theorem would compare the chamber volume of a Catalan-extreme sign word with that of ordinary sign words and then prove a lower bound for the chamber containing the actual prime-log vector.

\section{Three possible selector mechanisms}
\label{sec:selector_mechanisms}

\subsection{Hyperplane-chamber entropy}
\label{subsec:chamber_entropy}

The first route uses Definition~\ref{def:selector_chamber}.  For a finite excursion involving a finite collection of subset energies, let $\Omega_r$ be the region of positive weight space that preserves the relevant ordering and hence the same parity word.  Define a geometric selector energy
\begin{equation}
\mathcal{G}_r
=
-\log \nu(\Omega_r),
\label{eq:geometric_selector_energy}
\end{equation}
where $\nu$ is a normalized reference measure on an appropriate projective section of positive weight space.

\begin{question}
\label{question:chamber_vs_catalan_height}
Does a high Catalan excursion necessarily correspond to a small union of ordering chambers, with a bound of the form
\begin{equation}
\nu\{w:\text{induced excursion has height}\ge H\}
\le
\exp\!\left(-\Psi(H,L)\right),
\label{eq:chamber_height_rate}
\end{equation}
where $\Psi(H,L)$ is at least comparable to $H^2/L$ in the relevant regime?
\end{question}

If so, the remaining arithmetic problem becomes a Diophantine non-concentration statement for the prime-log vector relative to those chambers.  Proposition~\ref{prop:subset_gap} gives a first, very weak, deterministic separation from the collision walls.

\subsection{Shadow inequalities and parity restoration}
\label{subsec:shadow_restoration}

Bj\"orner obtains number-theoretic inequalities from combinatorial shadow functions associated with the shifted complex \cite{Bjorner2011}.  A family of many high-dimensional faces forces a sufficiently large lower-dimensional shadow.  Since adjacent dimensions contribute opposite signs to Euler characteristic, this is a plausible source of parity compensation.

\begin{warning}
\label{warn:shadow_temporal_gap}
A shadow inequality is primarily static: it says that certain lower-dimensional faces exist when a higher-dimensional face exists.  The Mertens problem is temporal along the filtration.  Many shadow faces may have entered much earlier.  The useful statement must therefore be a filtered shadow inequality controlling how much opposite parity is forced to enter within a bounded future energy interval.
\end{warning}

The corresponding research target is a dynamic inequality of the schematic form
\begin{equation}
\text{current parity imbalance }H
\quad\Longrightarrow\quad
\text{future opposite-parity influx}\ge \Phi(H,L),
\label{eq:dynamic_shadow_target}
\end{equation}
with a scale strong enough to feed into the Catalan energy $H^2/L$.

\subsection{Persistent topology}
\label{subsec:persistent_topology}

Equation \eqref{eq:mertens_euler_characteristic} can be refined by the Euler--Poincar\'e formula.  Writing $\beta_k(\Delta_x)$ for reduced Betti numbers,
\begin{equation}
\Mert(x)
=
\sum_{k\ge0}(-1)^{k-1}\beta_k(\Delta_x).
\label{eq:mertens_betti_alternation}
\end{equation}
Björner studied the asymptotic behaviour of these Betti numbers and proved that $\Delta_x$ has the homotopy type of a wedge of spheres \cite{Bjorner2011}.  A large Mertens excursion is therefore a sustained imbalance between even- and odd-dimensional homological contributions.

The filtration suggests studying not only $\beta_k(\Delta_x)$ at a fixed $x$ but the birth and death structure across $x$.  Persistent-homology constraints could in principle control the duration for which one parity dominates.  This route preserves information discarded by the scalar Euler characteristic.

\begin{question}
\label{question:persistent_parity}
Can one prove an interlacing, pairing, or energy inequality between persistent homology classes of adjacent parity whose projection onto Euler characteristic yields a bound on Catalan excursion height?
\end{question}

\section{Structural viability of the physical routes}
\label{sec:not_vacuous}

The physical structures above do not by themselves prove a new estimate for $\Mert$.  They are nevertheless mathematically relevant for four specific reasons.

First, the Catalan state space is exact, not fitted.  Proposition~\ref{prop:primitive_catalan_cardinality} follows from the deterministic zero-return geometry of the compressed Mertens walk.

Second, the physical rate $H^2/L$ is already visible in exact combinatorics.  Proposition~\ref{prop:meander_gaussian_tail} supplies a finite discrete Gaussian tail, while Brownian action and large-deviation theory explain the same scale in the continuum.

Third, the arithmetic selector has additional exact structure.  Equations \eqref{eq:mertens_euler_characteristic} and \eqref{eq:subset_sum_energy} show that the path is generated by a shifted simplicial filtration and an additive prime-log ordering.  Therefore the selector problem is not an arbitrary deterministic choice from Catalan space.

Fourth, Theorem~\ref{thm:selection_energy_implies_rh_scale} converts a concrete selector-energy inequality into square-root control.  This isolates a finite combinatorial target that can be attacked by geometry, topology, spectral theory, or statistical mechanics.

\section{Potential obstructions}
\label{sec:potential_obstructions}

\subsection{The selector may remain unconstrained at the needed scale}
\label{subsec:obstruction_selector}

The strongest obstruction is that the known simplicial and subset-sum structure may fail to prevent the prime-log ordering from entering Catalan sectors with very large selection energy.  Establishing such a prevention is the substantive new theorem required by the program.

\subsection{Static topology may not control filtered dynamics}
\label{subsec:obstruction_static_dynamic}

Euler characteristic and Betti-number identities can exhibit large cancellation only after many filtration steps.  RH-scale control requires cancellation at every sufficiently large scale.  Any topological approach must therefore retain filtration time and cannot rely solely on the terminal homotopy type.

\subsection{Parity barrier}
\label{subsec:obstruction_parity}

The selector records the parity $(-1)^{|S|}$ of prime subsets.  Many coarse counting and sieve mechanisms are weak precisely when asked to distinguish even from odd numbers of prime factors.  In the present language, a successful selector theorem must see parity at the level of the ordered subset-sum geometry rather than only the density of faces.

\subsection{Canonical collapse to a classical Dirichlet object}
\label{subsec:obstruction_canonical_collapse}

A canonical partition function built directly from the signed energies \eqref{eq:subset_sum_energy} gives
\begin{equation}
\sum_{S}(-1)^{|S|}e^{-sE(S)}
=
\prod_p(1-p^{-s})
=
\frac{1}{\zeta(s)}
\label{eq:canonical_collapse_zeta}
\end{equation}
for $\Re s>1$.  Equation \eqref{eq:canonical_collapse_zeta} is exact, but it collapses the excursion geometry into the familiar Euler product.  The new information sought here lies in pathwise objects---return times, Catalan shapes, heights, areas, chamber geometry, and filtered topology---before such a transform removes the temporal structure.

\begin{remark}
\label{rem:zeta_role}
Equation \eqref{eq:canonical_collapse_zeta} is used as a diagnostic boundary: a proposed physical construction that immediately reduces to this canonical partition function has returned to the classical analytic representation and has not exploited the Catalan excursion structure.
\end{remark}

\section{A hierarchy of possible sufficient principles}
\label{sec:sufficient_principles}

The following principles are ordered from closest to the desired estimate to more structural formulations.

\begin{conjecture}[Open Catalan selection principle]
\label{conj:open_catalan_selection}
For the compressed Mertens walk,
\begin{equation}
\Icat(m)=O(\log m).
\label{eq:open_catalan_selection_conjecture}
\end{equation}
\end{conjecture}

By Theorem~\ref{thm:selection_energy_implies_rh_scale}, Conjecture~\ref{conj:open_catalan_selection} implies \eqref{eq:mertens_sqrtlog_bound}.

\begin{conjecture}[Chamber non-concentration principle]
\label{conj:chamber_nonconcentration}
For a natural finite-dimensional normalization of the prime-log selector, the chamber containing the arithmetic weight vector has geometric energy at most logarithmic in the cutoff, and Catalan-extreme endpoint sectors have chamber energy bounded below by a positive multiple of $H^2/L$.
\end{conjecture}

\begin{conjecture}[Filtered parity-restoration principle]
\label{conj:filtered_parity_restoration}
There is a deterministic filtered shadow or persistence functional $\mathcal R$ such that a parity imbalance $H$ sustained for $L$ compressed steps forces
\begin{equation}
\mathcal R
\ge
c\frac{H^2}{L}-O(\log L),
\label{eq:parity_restoration_energy}
\end{equation}
and the arithmetic filtration satisfies $\mathcal R=O(\log m)$ up to compressed time $m$.
\end{conjecture}

These conjectures are not independent claims of truth.  They are candidate formulations of the same missing selector mechanism at different levels: direct Catalan volume, prime-log chamber geometry, and filtered topology.

\section{Research program}
\label{sec:research_program}

The route can be advanced without assuming a probabilistic model for M\"obius by attacking finite, exact questions.

\subsection{Exact finite combinatorics}
\label{subsec:program_exact}

For each actual excursion or meander prefix, compute
\begin{equation}
(L,H,A,S_{\mathrm{Cat}},\Icat),
\label{eq:finite_observables}
\end{equation}
where $L$ is duration, $H$ is maximum height or current endpoint, $A$ is area, $S_{\mathrm{Cat}}$ is Catalan entropy, and $\Icat$ is the open-sector selection energy.  The formulas for all ambient counts are exact or can be evaluated by stable dynamic programming/transfer matrices.

\subsection{Finite selector chambers}
\label{subsec:program_chambers}

For a finite set of primes and squarefree products, build the hyperplane arrangement \eqref{eq:subset_order_hyperplane}.  Determine which chambers generate Catalan excursions of a prescribed length and height.  The first structural question is whether the total chamber volume of height-$H$ words exhibits the same $H^2/L$ rate as the uniform Catalan ensemble or a stronger one.

\subsection{Dynamic shadows}
\label{subsec:program_shadows}

Translate Bj\"orner's static shadow inequalities into statements on increments of face numbers over multiplicative windows
\begin{equation}
[x,xe^{\delta}].
\label{eq:multiplicative_window}
\end{equation}
The target is a lower bound on opposite-parity arrivals conditional on a current Euler-characteristic imbalance.

\subsection{Transfer operators}
\label{subsec:program_transfer}

Construct a transfer operator whose states retain both Catalan height and a finite amount of arithmetic filtration data.  The ordinary Dyck transfer matrix has spectral gap scale $h^{-2}$; an arithmetic deformation could reveal whether the prime-log selector adds a confining or restoring term.

\subsection{Hydrodynamic coarse graining}
\label{subsec:program_hydrodynamic}

If a local transport rule emerges, derive its conservation law and large-deviation functional.  The TASEP example shows that Brownian-excursion fluctuations can coexist with Burgers hydrodynamics \cite{DerridaEnaudLebowitz2004}.  The corresponding Mertens question is whether the arithmetic selector generates an entropy condition that suppresses macroscopic parity shocks.

\section{Relation to existing literature}
\label{sec:literature_position}

The closest direct probabilistic work on M\"obius located in this literature survey is the study of Mussardo and LeClair, which reorganizes the nonzero M\"obius coefficients into a $\pm1$ walk and investigates Brownian/random-walk behaviour through statistical tests and block structures \cite{MussardoLeClair2021}.  The present program differs in its primary invariant: the exact Catalan state space attached to each return excursion and its open meander sectors.

On the mathematical-probability side, conditioned random walks converging to Brownian excursion or meander are classical \cite{Kaigh1976,DurrettIglehartMiller1977,CsakiMohanty1981}; the maximum of a random Dyck path and its exponential moments and deviations are quantitatively understood \cite{KhorunzhiyMarckert2009,Kousha2012}.  On the statistical-mechanics side, Dyck paths occur as vesicle/polymer configurations with genuine partition functions and free-energy differences \cite{OwczarekPrellberg2014,HaugPrellberg2015}.  On the quantum side, Fredkin-type spin chains realize Dyck-path superpositions as physical states \cite{DellAnnaEtAl2016,SalbergerKorepin2017}.  On the non-equilibrium side, Brownian excursions arise in exact TASEP fluctuation calculations, with Burgers-type hydrodynamics at macroscopic scale \cite{DerridaEnaudLebowitz2004}.

On the arithmetic-topological side, Bj\"orner's complex gives the exact identity \eqref{eq:mertens_euler_characteristic}, shifted-complex structure, Betti asymptotics, and shadow inequalities, while explicitly leaving the growth problem for $\Mert$ unresolved \cite{Bjorner2011}.  The potential new intersection is therefore
\begin{equation}
\boxed{
\text{filtered arithmetic selector}
+
\text{Catalan excursion state space}
+
\text{physical free-energy geometry}
}.
\label{eq:literature_intersection}
\end{equation}
The notebook treats the selector theorem as the decisive missing component rather than as an implicit randomness assumption.

\section{Conclusions and open questions}
\label{sec:conclusions}

The compressed Mertens function is an exact nearest-neighbour walk, and every complete zero-to-zero excursion is a primitive Dyck path counted by a Catalan number.  This yields a canonical finite state space for each excursion and a natural open-boundary meander space for every point inside an excursion.  Standard Catalan probability then supplies a quantitative geometric energy: large endpoint displacement or height is exponentially expensive on the scale $H^2/L$.

The central unresolved issue is selection.  The arithmetic path is generated by a deterministic ordering of prime-subset energies and simultaneously by a filtered shifted simplicial complex whose Euler characteristic is Mertens.  The project becomes viable only if one of these structures can be shown to bound the Catalan selection energy.

The most concrete first targets are:
\begin{enumerate}[label=(\roman*)]
\item prove nontrivial chamber-volume inequalities for parity words induced by positive additive subset weights;
\item derive filtered versions of Bj\"orner's shadow inequalities over multiplicative windows;
\item test whether persistent homology forces adjacent-parity compensation on the timescale of a Catalan excursion;
\item construct an arithmetic deformation of the Dyck transfer matrix and identify its spectral energy;
\item determine whether a current/transport formulation leads to an entropy-production principle analogous to non-equilibrium lattice gases.
\end{enumerate}

A successful result need not identify M\"obius with a random walk.  The decisive statement would instead be a deterministic restriction showing that the prime-generated selector cannot occupy Catalan sectors of superlogarithmic free energy.  Theorem~\ref{thm:selection_energy_implies_rh_scale} records one precise form in which such a statement would force square-root control.

\begin{thebibliography}{99}

\bibitem{Bjorner2011}
A.~Bj\"orner,
\emph{A cell complex in number theory},
Adv. Appl. Math. \textbf{46} (2011), 71--85.
DOI: 10.1016/j.aam.2010.09.007.
arXiv:1101.5704.

\bibitem{Kaigh1976}
W.~D.~Kaigh,
\emph{An invariance principle for random walk conditioned by a late return to zero},
Ann. Probab. \textbf{4} (1976), 115--121.
DOI: 10.1214/aop/1176996189.

\bibitem{DurrettIglehartMiller1977}
R.~T.~Durrett, D.~L.~Iglehart, and D.~R.~Miller,
\emph{Weak convergence to Brownian meander and Brownian excursion},
Ann. Probab. \textbf{5} (1977), 117--129.
DOI: 10.1214/aop/1176995895.

\bibitem{CsakiMohanty1981}
E.~Cs\'aki and S.~G.~Mohanty,
\emph{Excursion and meander in random walk},
Canad. J. Statist. \textbf{9} (1981), 57--70.
DOI: 10.2307/3315296.

\bibitem{KhorunzhiyMarckert2009}
O.~Khorunzhiy and J.-F.~Marckert,
\emph{Uniform bounds for exponential moment of maximum of a Dyck path},
Electron. Commun. Probab. \textbf{14} (2009), 327--333.
DOI: 10.1214/ECP.v14-1486.
arXiv:0902.1229.

\bibitem{Kousha2012}
T.~Kousha,
\emph{Asymptotic behavior and the moderate deviation principle for the maximum of a Dyck path},
Statist. Probab. Lett. \textbf{82} (2012), 340--347.
DOI: 10.1016/j.spl.2011.10.015.
arXiv:1008.0606.

\bibitem{OwczarekPrellberg2014}
A.~L.~Owczarek and T.~Prellberg,
\emph{Pressure exerted by a vesicle on a surface},
J. Phys. A: Math. Theor. \textbf{47} (2014), 215001.
DOI: 10.1088/1751-8113/47/21/215001.
arXiv:1311.2174.

\bibitem{HaugPrellberg2015}
N.~Haug and T.~Prellberg,
\emph{Uniform asymptotics of area-weighted Dyck paths},
J. Math. Phys. \textbf{56} (2015), 043301.
DOI: 10.1063/1.4917052.

\bibitem{DerridaEnaudLebowitz2004}
B.~Derrida, C.~Enaud, and J.~L.~Lebowitz,
\emph{The asymmetric exclusion process and Brownian excursions},
J. Stat. Phys. \textbf{115} (2004), 365--382.
DOI: 10.1023/B:JOSS.0000019833.35328.b4.
arXiv:cond-mat/0306078.

\bibitem{DellAnnaEtAl2016}
L.~Dell'Anna, O.~Salberger, L.~Barbiero, A.~Trombettoni, and V.~E.~Korepin,
\emph{Violation of cluster decomposition and absence of light cones in local integer and half-integer spin chains},
Phys. Rev. B \textbf{94} (2016), 155140.
DOI: 10.1103/PhysRevB.94.155140.
arXiv:1604.08281.

\bibitem{SalbergerKorepin2017}
O.~Salberger and V.~Korepin,
\emph{Entangled spin chain},
Rev. Math. Phys. \textbf{29} (2017), 1750031.
DOI: 10.1142/S0129055X17500313.

\bibitem{MussardoLeClair2021}
G.~Mussardo and A.~LeClair,
\emph{Randomness of M\"obius coefficients and Brownian motion: growth of the Mertens function and the Riemann hypothesis},
J. Stat. Mech. Theory Exp. (2021), 113106.
DOI: 10.1088/1742-5468/ac22fb.
arXiv:2101.10336.

\end{thebibliography}

\end{document}
